\documentclass[11pt,letterpaper]{article}

% --------------------------------------------------------------------------
% MOLE mid-semester design report template
% Replace text in angle brackets and the example rows in tables as the project
% develops. The document is intentionally structured to match the assignment
% brief; sections may be expanded or moved as needed.
% --------------------------------------------------------------------------

\usepackage[margin=1in]{geometry}
\usepackage{setspace}
\usepackage{graphicx}
\usepackage{booktabs}
\usepackage{longtable}
\usepackage{tabularx}
\usepackage{array}
\usepackage{multirow}
\usepackage{enumitem}
\usepackage{float}
\usepackage{pdflscape}
\usepackage{caption}
\usepackage{subcaption}
\usepackage{amsmath}
\usepackage{siunitx}
\usepackage{xcolor}
\usepackage{fancyhdr}
\usepackage{titlesec}
\usepackage{hyperref}
\usepackage{url}

\hypersetup{
    colorlinks=true,
    linkcolor=blue!50!black,
    citecolor=blue!50!black,
    urlcolor=blue!60!black,
    pdftitle={M.O.L.E. Design Report},
    pdfauthor={Group 3}
}

\setstretch{1.15}
\setlength{\parindent}{0pt}
\setlength{\parskip}{0.55em}
\setlist{itemsep=0.2em, topsep=0.25em}
\setcounter{tocdepth}{2}
\renewcommand{\arraystretch}{1.2}
\newcolumntype{Y}{>{\raggedright\arraybackslash}X}

\pagestyle{fancy}
\fancyhf{}
\lhead{MOLE Design Report}
\rhead{\thepage}
\setlength{\headheight}{14pt}
\renewcommand{\headrulewidth}{0.4pt}

\titleformat{\section}{\large\bfseries}{\thesection}{0.75em}{}
\titleformat{\subsection}{\normalsize\bfseries}{\thesubsection}{0.75em}{}

% ---------- Project-specific placeholders ----------
\newcommand{\projecttitle}{Multirotor-Operated Lab Experiments (MOLE)}
\newcommand{\coursetitle}{24671 Electromechanical Systems Design}
\newcommand{\groupname}{Group 3: M.O.L.E}
\newcommand{\reportdate}{<Month Day, Year>}
\newcommand{\placeholder}[1]{\textcolor{gray}{\textit{[#1]}}}

\newcommand{\insertfigure}[4][0.85\textwidth]{%
    \begin{figure}[H]
        \centering
        \IfFileExists{#2}{\includegraphics[width=#1]{#2}}{%
            \fbox{\parbox[c][2.25in][c]{#1}{\centering\placeholder{Insert figure: #2}}}}
        \caption{#3}
        \label{#4}
    \end{figure}
}

\begin{document}

% ============================== TITLE PAGE ================================
\begin{titlepage}
    \centering
    \vspace*{0.35in}
    \includegraphics[width=1.25in]{images/Carnegie_Mellon_University_seal.pdf}

    \vspace{0.45in}
    {\Large \coursetitle\par}
    \vspace{0.35in}
    {\LARGE\bfseries \projecttitle\par}
    \vspace{0.2in}
    {\Large Midsemester Design Report\par}
    \vfill

    \begin{tabular}{rl}
        \textbf{Date:} & \reportdate \\
        \textbf{Group:} & \groupname \\
    \end{tabular}

    \vspace{0.5in}
    \begin{tabular}{p{2.5in}p{2.5in}}
        \textbf{Team member} & \textbf{Signature} \\[0.35in]
        \placeholder{Name 1} & \hrulefill \\[0.35in]
        \placeholder{Name 2} & \hrulefill \\[0.35in]
        \placeholder{Name 3} & \hrulefill \\[0.35in]
        \placeholder{Name 4} & \hrulefill \\[0.35in]
        \placeholder{Name 5} & \hrulefill \\
    \end{tabular}
\end{titlepage}

\tableofcontents
\clearpage

% ========================== EXECUTIVE SUMMARY =============================
\section*{Executive Summary}
\addcontentsline{toc}{section}{Executive Summary}
\textbf{Problem.} \placeholder{In one paragraph, state the user problem, operating
context, and why it matters.}

\textbf{Product and conceptual design.} \placeholder{Summarize the proposed
system, its main features, and the design concept selected.}

\textbf{Final prototype.} \placeholder{State exactly what the final prototype
will demonstrate, and distinguish demonstrated functionality from deferred
features.}

\textbf{Results.} \placeholder{Summarize the prototype tests, the most important
quantitative results, and the principal design decisions that follow from them.}

\textbf{Takeaway.} \placeholder{Write this section last. It should stand alone as
an abstract of the full report.}

\clearpage

% =========================== PROBLEM DEFINITION ===========================
\section{Problem Definition}
\subsection{Problem Description}
\placeholder{Describe the problem in measurable terms. Explain the current
workflow, the users affected, the environment in which the system operates,
and the consequences of not solving the problem. Define any technical terms.}

\subsection{Markets Addressed}
\placeholder{Identify the market or application segments, expected users,
purchasers, and deployment environments. Explain the initial beachhead market
and possible future markets.}

\subsection{Primary Competitors and Existing Solutions}
\placeholder{Compare direct competitors, adjacent products, and the current
manual or improvised solution. Use consistent comparison criteria.}

\begin{table}[H]
\centering
\caption{Preliminary competitor comparison.}
\label{tab:competitors}
\begin{tabularx}{\textwidth}{lYYYY}
\toprule
\textbf{Competitor/solution} & \textbf{Core offering} & \textbf{Strengths} &
\textbf{Limitations} & \textbf{Web link} \\
\midrule
\placeholder{Competitor 1} & \placeholder{...} & \placeholder{...} &
\placeholder{...} & \url{https://example.com} \\
\placeholder{Competitor 2} & \placeholder{...} & \placeholder{...} &
\placeholder{...} & \url{https://example.com} \\
\placeholder{Current practice} & \placeholder{...} & \placeholder{...} &
\placeholder{...} & \placeholder{N/A} \\
\bottomrule
\end{tabularx}
\end{table}

\insertfigure[0.75\textwidth]{images/competitor-1.png}{Key competitor or
existing solution. Replace the filename and cite the source.}{fig:competitor}

\subsection{Assumptions and Constraints}
\begin{itemize}
    \item \placeholder{Technical assumption, e.g., available power, network,
    or lab infrastructure.}
    \item \placeholder{User or operational assumption.}
    \item \placeholder{Safety, budget, schedule, size, weight, or regulatory
    constraint.}
    \item \placeholder{What is explicitly out of scope for this project?}
\end{itemize}

% ===================== STAKEHOLDERS AND CUSTOMER NEEDS ====================
\section{Stakeholders and Customer Needs}
\subsection{Stakeholder Identification}
\placeholder{List primary users, customers, operators, maintainers, instructors,
administrators, suppliers, regulators, and other parties affected by the
design. Explain each stakeholder's influence and interest.}

\begin{table}[H]
\centering
\caption{Stakeholder map.}
\label{tab:stakeholders}
\begin{tabularx}{\textwidth}{lYYY}
\toprule
\textbf{Stakeholder} & \textbf{Role/context} & \textbf{Goals} &
\textbf{Influence/priority} \\
\midrule
\placeholder{Primary user} & \placeholder{...} & \placeholder{...} &
\placeholder{High} \\
\placeholder{Customer/purchaser} & \placeholder{...} & \placeholder{...} &
\placeholder{High} \\
\placeholder{Maintainer} & \placeholder{...} & \placeholder{...} &
\placeholder{Medium} \\
\bottomrule
\end{tabularx}
\end{table}

\subsection{External Outreach}
\placeholder{Describe interviews, surveys, observations, expert consultations,
or other outreach. State who participated, when, how many responses were
obtained, and how the evidence was analyzed.}

\subsection{Internal Need Generation}
\placeholder{Describe needs generated from team brainstorming, engineering
analysis, course requirements, and lessons from early prototypes.}

\subsection{Organization and Classification}
\placeholder{Explain how raw statements were converted into needs, normalized,
de-duplicated, and organized by user, function, or priority.}

\subsection{Final Customer Needs}
\begin{longtable}{p{0.11\textwidth}p{0.48\textwidth}p{0.16\textwidth}p{0.15\textwidth}}
\caption{Final customer needs.}\label{tab:customer-needs}\\
\toprule
\textbf{ID} & \textbf{Need statement} & \textbf{Source} & \textbf{Priority}\\
\midrule
\endfirsthead
\toprule
\textbf{ID} & \textbf{Need statement} & \textbf{Source} & \textbf{Priority}\\
\midrule
\endhead
\placeholder{N1} & \placeholder{The system should...} & \placeholder{Interview/survey} &
\placeholder{1--5} \\
\placeholder{N2} & \placeholder{The system should...} & \placeholder{Observation} &
\placeholder{1--5} \\
\placeholder{N3} & \placeholder{The system should...} & \placeholder{Team analysis} &
\placeholder{1--5} \\
\bottomrule
\end{longtable}

% =========================== TARGET SPECIFICATIONS =======================
\section{Target Specifications}
\subsection{Mapping Needs to Metrics}
\placeholder{Explain how each important need was translated into a measurable
engineering metric, including units, direction of improvement, and test method.}

\begin{table}[H]
\centering
\caption{Need-to-metric mapping.}
\label{tab:need-metric}
\begin{tabularx}{\textwidth}{lYYYY}
\toprule
\textbf{Need ID} & \textbf{Metric} & \textbf{Units} & \textbf{Why it matters} &
\textbf{Verification method} \\
\midrule
\placeholder{N1} & \placeholder{Response time} & \si{\second} &
\placeholder{...} & \placeholder{Bench test} \\
\placeholder{N2} & \placeholder{Accuracy} & \si{\percent} &
\placeholder{...} & \placeholder{Experiment} \\
\bottomrule
\end{tabularx}
\end{table}

\subsection{Competitive Analysis}
\placeholder{Compare competitor performance against the metrics above. Identify
gaps and opportunities rather than only listing features.}

\subsection{Target Specifications}
\begin{table}[H]
\centering
\caption{Final target specifications.}
\label{tab:specifications}
\begin{tabularx}{\textwidth}{lYYYY}
\toprule
\textbf{Metric} & \textbf{Minimum} & \textbf{Target} & \textbf{Ideal} &
\textbf{Test method} \\
\midrule
\placeholder{Metric 1} & \placeholder{...} & \placeholder{...} &
\placeholder{...} & \placeholder{...} \\
\placeholder{Metric 2} & \placeholder{...} & \placeholder{...} &
\placeholder{...} & \placeholder{...} \\
\placeholder{Metric 3} & \placeholder{...} & \placeholder{...} &
\placeholder{...} & \placeholder{...} \\
\bottomrule
\end{tabularx}
\end{table}

% ============================ CONCEPT GENERATION ==========================
\section{Concept Generation}
\subsection{Functional Decomposition}
\placeholder{Define the overall system function and decompose it into
subfunctions. A block diagram should show inputs, outputs, flows, and system
boundaries.}

\insertfigure{images/functional-decomposition.png}{Functional decomposition
of the proposed system.}{fig:functional-decomposition}

\subsection{Concept Generation Activities}
\subsubsection{External Search}
\placeholder{Describe literature, patent, product, standards, and web searches.
Record the sources and explain what design principles were adopted or rejected.}

\subsubsection{Internal Search}
\placeholder{Describe brainstorming, analogous products, morphological charts,
sketching, workshops, and other team activities.}

\subsection{Subsystem Concepts}
\placeholder{Present the alternatives generated for each major subsystem. Move
additional sketches to Appendix~\ref{app:concept-sketches}.}

\begin{table}[H]
\centering
\caption{Morphological chart of subsystem alternatives.}
\label{tab:morphological}
\begin{tabularx}{\textwidth}{lYYY}
\toprule
\textbf{Subsystem} & \textbf{Alternative A} & \textbf{Alternative B} &
\textbf{Alternative C} \\
\midrule
\placeholder{Subsystem 1} & \placeholder{...} & \placeholder{...} &
\placeholder{...} \\
\placeholder{Subsystem 2} & \placeholder{...} & \placeholder{...} &
\placeholder{...} \\
\placeholder{Subsystem 3} & \placeholder{...} & \placeholder{...} &
\placeholder{...} \\
\bottomrule
\end{tabularx}
\end{table}

\subsection{Full System Concepts}
\placeholder{Show the full-system combinations considered from the subsystem
alternatives. Give each concept a short name and describe its operating
principle, advantages, risks, and unresolved questions.}

\insertfigure{images/system-concepts.png}{Overview of full-system concepts
considered.}{fig:system-concepts}

% ============================ CONCEPT SELECTION ============================
\section{Concept Selection}
\subsection{Selection Process}
\placeholder{Describe screening, feasibility checks, a Pugh matrix, weighted
decision matrix, or other methods. Explain criteria, weights, scoring scale,
and how disagreements were resolved.}

\begin{table}[H]
\centering
\caption{Weighted concept-selection matrix.}
\label{tab:selection}
\begin{tabularx}{\textwidth}{lYYYY}
\toprule
\textbf{Criterion (weight)} & \textbf{Concept A} & \textbf{Concept B} &
\textbf{Concept C} & \textbf{Rationale} \\
\midrule
\placeholder{Criterion 1 (0.25)} & \placeholder{score} & \placeholder{score} &
\placeholder{score} & \placeholder{...} \\
\placeholder{Criterion 2 (0.25)} & \placeholder{score} & \placeholder{score} &
\placeholder{score} & \placeholder{...} \\
\placeholder{Criterion 3 (0.25)} & \placeholder{score} & \placeholder{score} &
\placeholder{score} & \placeholder{...} \\
\placeholder{Criterion 4 (0.25)} & \placeholder{score} & \placeholder{score} &
\placeholder{score} & \placeholder{...} \\
\midrule
\textbf{Weighted total} & \textbf{\placeholder{...}} & \textbf{\placeholder{...}} &
\textbf{\placeholder{...}} & \\
\bottomrule
\end{tabularx}
\end{table}

\subsection{Top Concept(s): Assessment and Justification}
\placeholder{Justify the selected concept using evidence, not only preference.
Discuss tradeoffs, risks, expected performance, cost, manufacturability, and
alignment with customer needs.}

\subsection{Detailed Model of Top Concept(s)}
\placeholder{Include CAD, architecture diagrams, circuit diagrams, software
models, interface descriptions, and key dimensions.}

\insertfigure{images/top-concept-cad.png}{Detailed model of the selected
concept.}{fig:top-concept}

% ============================ PROTOTYPE SUMMARY ============================
\section{Prototype Summary}
\subsection{Prototype Design}
\subsubsection{Questions to Answer}
\begin{enumerate}
    \item \placeholder{Question about technical feasibility.}
    \item \placeholder{Question about user interaction or workflow.}
    \item \placeholder{Question about performance against a target specification.}
    \item \placeholder{Question about risk, reliability, or safety.}
\end{enumerate}

\subsubsection{Design and Construction Process}
\placeholder{Describe the prototype architecture, materials, electronics,
software, fabrication, assembly, and integration. Explain the sequence of
decisions and document major iterations. Include photographs with captions.}

\insertfigure[0.48\textwidth]{images/prototype-overview.jpg}{Prototype
overview.}{fig:prototype-overview}
\insertfigure[0.48\textwidth]{images/prototype-build.jpg}{Prototype during
construction or integration.}{fig:prototype-build}

\subsection{Prototype Testing}
\subsubsection{Test Plan}
\placeholder{For each experiment, state the objective, setup, independent and
dependent variables, controls, sample size, procedure, and pass/fail criterion.}

\begin{table}[H]
\centering
\caption{Prototype test plan.}
\label{tab:test-plan}
\begin{tabularx}{\textwidth}{lYYYY}
\toprule
\textbf{Test} & \textbf{Question} & \textbf{Measured output} &
\textbf{Pass criterion} & \textbf{Status} \\
\midrule
\placeholder{T1} & \placeholder{...} & \placeholder{...} &
\placeholder{...} & \placeholder{Planned/complete} \\
\placeholder{T2} & \placeholder{...} & \placeholder{...} &
\placeholder{...} & \placeholder{Planned/complete} \\
\bottomrule
\end{tabularx}
\end{table}

\subsubsection{Data Analysis}
\placeholder{Present raw-data processing, uncertainty, summary statistics,
comparisons to targets, and limitations. Do not show only successful trials.}

\insertfigure{images/test-results.png}{Prototype test results. Replace with a
plot generated from the measured data.}{fig:test-results}

\begin{table}[H]
\centering
\caption{Summary of prototype test results.}
\label{tab:test-results}
\begin{tabularx}{\textwidth}{lYYY}
\toprule
\textbf{Metric} & \textbf{Target} & \textbf{Measured result} &
\textbf{Interpretation / next action} \\
\midrule
\placeholder{Metric 1} & \placeholder{...} & \placeholder{...} &
\placeholder{...} \\
\placeholder{Metric 2} & \placeholder{...} & \placeholder{...} &
\placeholder{...} \\
\bottomrule
\end{tabularx}
\end{table}

\subsubsection{Design Implications}
\placeholder{Explain how the test results changed the next design iteration,
specifications, concept choice, risk register, or test plan.}

% ============================ PROJECT PLANNING =============================
\section{Project Planning}
\subsection{Schedule and Gantt Chart}
\placeholder{Summarize completed work, current work, dependencies, milestones,
and schedule risks. Include a Gantt chart exported as a high-resolution image
or generated in this section.}

\insertfigure{images/gantt-chart.png}{Project Gantt chart.}{fig:gantt}

\subsection{Prototype Planning}
\subsubsection{Next Prototype}
\placeholder{Provide a detailed build plan: components, dimensions, interfaces,
fabrication steps, software work, responsible team member, cost, and expected
completion date.}

\subsubsection{Experiments}
\placeholder{List the experiments that will be conducted on the next prototype.
For each experiment, state the question, setup, data to collect, and decision
rule.}

\subsubsection{Expected Final Prototype Scope}
\placeholder{Describe the features expected in the final prototype. Clearly
separate must-have features from features deferred until a later product
version. Explain why features were deferred.}

\begin{table}[H]
\centering
\caption{Planned final-prototype features.}
\label{tab:final-features}
\begin{tabularx}{\textwidth}{lYYY}
\toprule
\textbf{Feature} & \textbf{Final prototype?} & \textbf{Evidence of completion} &
\textbf{Rationale / deferral plan} \\
\midrule
\placeholder{Feature 1} & \placeholder{Yes/No} & \placeholder{Test/demo} &
\placeholder{...} \\
\placeholder{Feature 2} & \placeholder{Yes/No} & \placeholder{Test/demo} &
\placeholder{...} \\
\bottomrule
\end{tabularx}
\end{table}

% ============================== REFERENCES ===============================
\clearpage
\begin{thebibliography}{99}
\bibitem{competitor1}
\placeholder{Author or organization.}
\emph{\placeholder{Title of source or competitor page}.}
\url{https://example.com}, \placeholder{publication date or access date}.

\bibitem{method}
\placeholder{Add course text, design-method reference, standard, paper, or
other source used in the report.}
\end{thebibliography}

% ============================== APPENDICES ================================
\appendix
\section{Individual Concept Sketches}
\label{app:concept-sketches}
\placeholder{Include at least one clearly labeled page or figure for each team
member's individual concepts.}

\insertfigure{images/sketch-member-1.png}{Concept sketches by
\placeholder{Team member 1}.}{fig:sketch-member-1}
\insertfigure{images/sketch-member-2.png}{Concept sketches by
\placeholder{Team member 2}.}{fig:sketch-member-2}

\section{Survey Instrument and Interview Questions}
\placeholder{Include the complete survey, interview script, consent language,
and response scale. Do not include personally identifying information unless
permission has been obtained.}

\begin{enumerate}
    \item \placeholder{Question 1}
    \item \placeholder{Question 2}
    \item \placeholder{Question 3}
\end{enumerate}

\section{Work Distribution}
\placeholder{Describe each team member's specific contribution to design
exercises, prototyping, testing, analysis, and report writing.}

\begin{longtable}{p{0.2\textwidth}p{0.36\textwidth}p{0.36\textwidth}}
\caption{Work distribution.}\label{tab:work-distribution}\\
\toprule
\textbf{Team member} & \textbf{Design/prototyping contribution} &
\textbf{Report contribution} \\
\midrule
\endfirsthead
\toprule
\textbf{Team member} & \textbf{Design/prototyping contribution} &
\textbf{Report contribution} \\
\midrule
\endhead
\placeholder{Name 1} & \placeholder{Specific tasks} & \placeholder{Sections,
figures, analysis, editing} \\
\placeholder{Name 2} & \placeholder{Specific tasks} & \placeholder{Sections,
figures, analysis, editing} \\
\placeholder{Name 3} & \placeholder{Specific tasks} & \placeholder{Sections,
figures, analysis, editing} \\
\placeholder{Name 4} & \placeholder{Specific tasks} & \placeholder{Sections,
figures, analysis, editing} \\
\bottomrule
\end{longtable}

\section{Additional Supporting Material}
\placeholder{Add detailed calculations, complete raw data, risk register,
additional concept sketches, CAD drawings, code excerpts, or other material
that supports the report but would interrupt the main narrative.}

\end{document}